\chapter{Поле проекторного параметра порядка и спектр его мод}

После вывода одноосной фазы необходимо перестать говорить о материи только
как о метафоре кристаллизации. Первым полем, которое уже действительно
следует из проекта, является пространственно зависящий параметр порядка
\begin{equation}
 Q(x)=R(x)-\frac{I_3}{3},
 \qquad Q^T=Q,
 \qquad \Tr Q=0.
 \label{eq:v6-projective-Q-field}
\end{equation}
Это вещественное симметричное бесследовое поле с пятью локальными
компонентами.

\section{Упорядоченный фон и вакуумная орбита}

При точке сосуществования
\begin{equation}
 \beta_c=1.542669540860\ldots
\end{equation}
предыдущий фазовый аудит получил спектр
\begin{equation}
 R_*=\operatorname{diag}(a_*,b_*,b_*),
 \qquad
 a_*=0.9121666003\ldots,
 \quad
 b_*=0.04391669984\ldots.
 \label{eq:v6-projective-ordered-background}
\end{equation}
Стабилизатор такого состояния есть $O(2)$, поэтому множество равных
минимумов равно
\begin{equation}
 \mathcal M_{\rm vac}=SO(3)/O(2)=\mathbb{RP}^2.
 \label{eq:v6-projective-vacuum-orbit}
\end{equation}
Прямой поворот фона на 101 точке орбиты сохранил энергию с остатком ниже
$6.2\cdot10^{-16}$; сто случайных ортогональных преобразований дали
остаток ниже $3.0\cdot10^{-15}$.

Тензорный параметр порядка и разложение на одноосные, двухосные и
директорные моды являются стандартным языком теории Ландау--де Жена; см.
DOI \path{10.1103/PhysRevE.55.4367}. Топологическая классификация дефектов
пространством параметра порядка опирается на DOI
\path{10.1103/RevModPhys.51.591}.

\section{Пять локальных полей}

В ортонормированном базисе симметричных бесследовых матриц гессиан полного
канонического функционала раскладывается на три сектора:
\begin{enumerate}
 \item одна амплитудная мода вдоль
 $\operatorname{diag}(2,-1,-1)/\sqrt6$;
 \item две двухосные моды, расщепляющие поперечную пару $b_*,b_*$;
 \item две директорные моды, поворачивающие выделенную ось.
\end{enumerate}
Численные кривизны равны
\begin{equation}
 \operatorname{Spec}\Hess_{R_*}\mathcal F
 =\{4.5081528\ldots,
 17.8845270\ldots,17.8845270\ldots,0,0\}.
 \label{eq:v6-projective-field-hessian}
\end{equation}

\begin{theorem}[Первичный бозонный спектр упорядоченной фазы]
Локальное поле $Q$ содержит одну массивную амплитудную моду, две
вырожденные массивные двухосные моды и две безмассовые директорные моды
Голдстоуна. Последние являются касательными к $\mathbb{RP}^2$-орбите и не
вставлены как отдельные поля.
\end{theorem}

Это уже физика коллективных возбуждений. Но директорные моды нельзя
автоматически называть калибровочными бозонами: для этого $SO(3)$ должно
быть локализовано, а его связность, кинетическая норма и физическая факторизация
должны быть выведены из родительской суперсвязности.

\section{Первые честные кандидаты на материю}

Для вакуумного пространства $\mathbb{RP}^2$ известны
\begin{equation}
 \pi_1(\mathbb{RP}^2)=\mathbb Z_2,
 \qquad
 \pi_2(\mathbb{RP}^2)=\mathbb Z,
 \qquad
 \pi_3(\mathbb{RP}^2)=\mathbb Z.
 \label{eq:v6-projective-homotopy-sectors}
\end{equation}
Поэтому поле $Q(x)$ допускает соответственно линейные дисклинации,
целочисленные точечные ежи и трёхмерные хопфовы текстуры. Именно эти
топологические сектора, а не малые гладкие колебания сами по себе, являются
первым строгим кандидатом на материю как устойчивый дефект среды.

Том V уже установил необходимые мосты, но не полное отождествление:
проективный ёж требует спинорного подъёма, хопфова линия переносит
ориентацию, а конструкция Ботта--Тёплица несёт Real-пару $+15/-15$.
Следовательно, текущий результат создаёт пространство полей и классы
материальных объектов, но ещё не доказывает, что конкретный класс есть
электрон, нейтрино или другая наблюдаемая частица.

\section{Следующий различающий тест}

Следующая проверка должна вернуть пространственные производные и вычислить для
каждого сектора:
\begin{enumerate}
 \item существование конечной энергии;
 \item число мод трансляции, вращения и изменения размера;
 \item баланс квадратичной и скирмовской жёсткости;
 \item спектр локализованных фермионных нулевых мод;
 \item совместимость со меткой спинорного подъёма $+15/-15$.
\end{enumerate}

\begin{nogocriterion}[Топологический сектор ещё не является частицей]
Ненулевая группа гомотопий доказывает наличие класса конфигураций, но не
конечную массу, устойчивый радиус, квантовую статистику или соответствие
наблюдаемой частице.
\end{nogocriterion}

\begin{protocol}
\begin{enumerate}
 \item первое выведенное бозонное поле: \textbf{$Q=R-I_3/3$};
 \item число локальных компонент: \textbf{пять};
 \item массивные моды: \textbf{одна амплитудная и две двухосные};
 \item моды Голдстоуна: \textbf{две директорные};
 \item вакуумная орбита: \textbf{$\mathbb{RP}^2$};
 \item линейные, точечные и хопфовы сектора: \textbf{существуют};
 \item конечная энергия каждого сектора: \textbf{ещё не доказана};
 \item отождествление с наблюдаемыми частицами: \textbf{не сделано};
 \item следующая проверка:
 \textbf{пространственная энергия и спектр проекторных дефектов};
 \item переход к материи и полям: \textbf{начат, но не завершён}.
\end{enumerate}
\end{protocol}

Машинный сертификат сохранён в
\path{s2t_v6_projective_order_parameter_field_spectrum_gate_results.json}.